\documentclass[12pt]{article}
\usepackage[utf8]{inputenc}
\usepackage[T1]{fontenc}
\usepackage{hyperref}
\title{Wejewa Language Resources for K-12 Curriculum Development}
\author{Compiled by Claude Sonnet 5.5 (claude-sonnet-5-5) for Andr\'as Kornai}
\date{}
\begin{document}
\maketitle

\section{Introduction}
Wejewa (also Weyewa or Wewewa; ISO 639-3 wew) is an Austronesian language of Sumba, Indonesia. Sources cite roughly 55,000 to 65,000 speakers. Search turned up no full grammar, no comprehensive dictionary and no school materials. The resources below are partial.

\section{Primary Educational Resources}
No primers, textbooks or school curricula in Wejewa were found in search results.

\section{Linguistic and Reference Materials}
\begin{itemize}
\item Wikipedia article on the Wejewa language, with a consonant and vowel inventory. It cites a 2018 Indonesian-language thesis on the phonology of the Wewewa language in the village of Weekombaka. \url{https://en.wikipedia.org/wiki/Wejewa_language}
\item Verbix Wejewa verb list. This is a short list with English glosses, not a dictionary. \url{https://docs.verbix.com/Languages/Wejewa}
\item Omniglot page on Wejewa, covering writing and orthography history. \url{https://omniglot.com/writing/wejewa.htm}
\item Wiktionary entry for Weyewa. \url{https://en.wiktionary.com/wiki/Weyewa}
\end{itemize}

\section{Religious and Community Resources}
Omniglot reports that the language was first written down in the 1930s by Dr. Onvlee and that a New Testament translation was published in 1970. No URL for the scripture text itself was seen in search results.

\section{Digital and Audio Resources}
\begin{itemize}
\item PARADISEC collection WEW2018, with recordings of texts and elicitation sessions (tense, aspect and pronouns). Access is subject to PARADISEC conditions. \url{https://catalog.paradisec.org.au/repository/WEW2018}
\item An example item from that collection. \url{https://catalog.paradisec.org.au/collections/WEW2018/items/MON022}
\item Research Data Australia record, ``Recordings of various texts, Tenggara Indonesia''. \url{https://researchdata.edu.au/recordings-various-texts-tenggara-indonesia/1426096}
\end{itemize}

\section{Curriculum Development Context}
With about 55,000 to 65,000 speakers, the language is stable, but the search found no evidence that it is taught in schools. A curriculum would need to build its own primers and reading materials. The PARADISEC recordings and the Bible translation are possible starting points for this.

\end{document}